\section{Motivation}
Nucelar energy supplied about 20\% of all energy and abot 54\% 
of all carbon-free energy in the US in 2021 \hl{citation}, making nuclear 
energy the third largest producer of energy and the largest 
producer of carbon-free energy in the US. Because of its role 
as a carbon free supply of base-load power, nuclear energy is expected 
to play a role in meeting carbon emission goals. The current nuclear energy 
indsutry is comprised of two reactor designs: \glspl{PWR} and \glspl{BWR}, 
which both fall under the classification of \glspl{LWR}.
Both designs use uranium dioxide pellets as fuel, with the uranium 
enriched to no more than 5\% uranium-235. These types of reactors 
have commercially operated in the US since 1957. The \glspl{LWR} 
currently deployed will be retiring over the next 30 years, with the 
last license expiring in 2055 \cite{nei_us_2021}. Therefore, if 
nuclear energy is to continue to contribute to the current US energy 
grid and assit in meeting carbon emission goals \hl{citation}, new 
reactors must be deployed. 

New \glspl{LWR} are being deployed around the world \hl{Citation}, but 
new reactor designs, often called advanced reactors, are being developed 
\cite{hussain_advances_2018}to 
replace the current fleet or reactors or expand it. These advanced reactors 
cover a large swath of the design space, from their energy output to their 
fuel form, to their cycle length. Some of these new designs are being 
investigated for how they could be built and contribute to the US energy 
grid. Much of these investigations are through the \gls{DOE} \gls{ARDP}
\hl{citation}, which aims to support research and the building of 
demostrations for a select set of the new designs. Some of the goals 
of the program are to \hl{find more}. 

One central question that has arisen from the \gls{ARDP} is the development 
of supply chains to supprot the nuclear fuel cycle of the advanced reactors 
\hl{ciation}. One key difference in the fuel for these advanced 
reactor designs and the fuel for the current \glspl{LWR} is the enrichment 
level of the fuel. Current \glspl{LWR} require fuel enriched to less than 
5\% uranium-235, but the advanced reactor designs require fuel enriched to 
up to 20\% uranium-235. There is currently no commercial supply of 
uranium enriched to this level in the US. Therefore the \gls{DOE} is 
investigating how to develop a supply chain for uranium enriched to this 
higher level. 

The current nucler fuel cycle in the US is based on supplying reactors 
with fuel comprised of enriched uranium. Natural uranium consists of 
about 0.7\% uranium-235 by weight, with most of the rest of the weight being 
uranium-238. Enriched uranium sent to reactors is enriched up to 4.95\% 
uranium-235 \hl{Find citation, NRC regs}. The US nuclear industry is 
interested in developing a supply chain for \gls{HALEU}, which is 
uranium enriched to between 5-20\% uranium-235 \hl{Find citation}. \gls{HALEU}
has a variety of potential uses, including for medical isotope production,
test reactors, and thermal propulsion for space exploration \cite{nagley_ha-leu_2020}.
For nuclear power reactors, \gls{HALEU} fuel is expected to provide benefts 
over current fuel enrichment levels, such as higher burnups \hl{citation},
longer cycle times, and 
reducing the \gls{LCOE} for the reactors \cite{carlson_economic_2020}. These 
benefits have led many advanced reactor designs to require 
\gls{HALEU} fuel instead of fuel enriched to less than 5\% uranium-235. 
Many of the current \glspl{LWR} have license expiration dates before 2050, 
so it is very likely that advanced reactors will be built in order to 
provide energy and help countries meet their carbon emission goals \hl{citation}.

Two methods can be used to produce \gls{HALEU}: enrich natural uranium to the 
required level or downblend \gls{HEU} to the required level. The current US
nuclear commercial fuel cycle does not have the capability to produce 
\gls{HALEU} using either of these methods \hl{citation}. Enriching uranium 
to the required enrichment level is limited by the amount of natural uranium 
that can be mined and the \gls{SWU} capacity of the available facilities. 
Downblending \gls{HEU} is limited by the amount of \gls{HEU} avaialable and 
the downblending capacity of the available facilities. 

There is currently no commercial facility in the US that can enrich 
uranium to produce \gls{HALEU} \cite{hussain_nei_2018}, which has led 
many to consider downblending \gls{HEU} to obtain an initial stockpile of 
\gls{HALEU}. 
Currently, there is only one facility in the US commercially licensed to 
downblend \gls{HEU}, the BWXT Nuclear Fuel Services Inc. facility in 
Erwin, TN, which is expected to have the capacity to downblend 1-2 
MT of \gls{HEU} into up to 10 MT of \gls{HALEU} each year \cite{nagley_ha-leu_2020}.
The \gls{HEU} being considered for downblending mostly comes from the spent 
fuel of the \gls{EBR} reactor at \gls{INL} \cite{patterson_haleu_2019}. From 
this stockpile, there is about three metric tonnes of \gls{HEU} that 
could be downblended \cite{patterson_haleu_2019}.

The context of this PhD dissertation is to investigate the impacts  
of deploying reactor deisgns that require \gls{HALEU} for 
fuel. This investigation is performed through modeling the transition from 
the current fleet of US \glspl{LWR} to \gls{HALEU}-fueled advanced reactors, 
then performs sensitivity analysis an optimization of the transitions. 
Finally, the investigation concludes by modeling the neutonics of select 
\gls{HALEU}-fueled reactors with impure fuel from downblending. 

\section{The nuclear fuel cycle}
The nuclear fuel cycle encompasses the activities and processes 
for the use of fissile materials in fission nuclear reactors \cite{tsoulfanidis_nuclear_2013}. 
Uranuim, thorium, or even plutonium can be used as the 
fuel for the reactor, with uranium being the most widely used fuel for 
power reactors \hl{FIND CITATION}. This work focuses on a uranium-based 
fuel cycle. The nuclear fuel cycle begins with the mining of uranium ores from the earth 
and ends with either the final disposal of the radioactive waste produced 
\cite{tsoulfanidis_nuclear_2013}. A fuel cycle can vary based on the
design of the reactor(s) deployed, such as if enrichment is required.  
Two possible variations on the nuclear fuel cycle 
include the once-through and recycling fuel cycles 
\cite{tsoulfanidis_nuclear_2013}. Figure \ref{fig:fuel_cycle} shows 
an example of all possible steps in a nuclear fuel cycle. 

\begin{figure}
    \centering
    \includegraphics[scale=0.8]{nuclear_fuel_cycle.jpg}
    \caption{A complete view of options available for a once-through and 
    a recycling nuclear 
    fuel cycle. Graphic from \protect\cite{noauthor_stages_2020}.}
    \label{fig:fuel_cycle}
\end{figure}

\subsection{Once-through fuel cycle}
A once-through fuel cycle is characterized by fuel going through one 
pass through a reactor before being disposed of. In Figure \ref{fig:fuel_cycle},
this fuel cycle option is shown through the gray line from the storage 
step to the disposal step. None of the material flows from the Reprocessing 
Facility in Figure \ref{fig:fuel_cycle} are included in a once-through 
fuel cycle. This is the fuel cycle option currently employed in the 
United States. 
%\begin{figure}
    \centering
    \begin{tikzpicture}[node distance=1.5cm]
        \node (ore) [facility] {Uranium ore};
        \node (enrichment) [facility, right of=ore, xshift=1cm]{Enrich};
        \node (llwsink) [facility, below of=enrichment]{Depleted uranium};
        \node (fabrication) [facility, right of=enrichment, xshift=1cm]{Fuel};
        \node (reactor) [facility, above of=fabrication, xshift=2cm]{Power Plant};
        \node (spentfuel) [facility, below of=reactor, xshift=2cm]{Interim Storage};
        \node (sinkhlw) [facility, below of=spentfuel, xshift=2cm]{Geologic Repository};

        \draw [arrow] (ore) --  (enrichment); 
        \draw [arrow] (enrichment) -- (fabrication);
        \draw [arrow] (enrichment) -- (llwsink);
        \draw [arrow] (fabrication) |- (reactor);
        \draw [arrow] (reactor) -| (spentfuel);
        \draw [arrow] (spentfuel) -| (sinkhlw);

        \end{tikzpicture}
    \caption{Material flow for a once-through nuclear fuel cycle. All material is 
    disposed of after use in a reactor. Adapted from \protect\cite{wigeland_identification_2011}.}
    \label{fig:once_through_background}
\end{figure}

\subsection{Recycling fuel cycle}
A recycling fuel cycle is characterized by fuel going through additional 
passes through a reactor after a separtation step to remove fission 
products. There are a variety of methods to employ recycling in a 
fuel cycle. One distinction between recycling options is the number of 
times fuel is reprocessed. If fuel is reprocessed and burned in a reactor
for a finite number of times before being disposed of, then this recycling 
scheme is often referred to as a limited recycle. However, if the fuel is 
repeatedly reprocessed and burned in a reactor, then this is 
often referred to as a continuous reprocessing scheme. 

There are two primary methods to reprocess \gls{SNF}: aqueous reprocessing 
and pyroprocessing. Aqueous reprocessing relies on redox reactions of 
select actinides (typically uranium and plutonium) to remove the 
actinides from an aqueous phase to an organic phase, while the non-desired 
materials in the \gls{SNF} remain in the aqueous phase \hl{citation}. There 
are a variety of methods to perform aqueous reprocessing, but the primary 
method is the PUREX process. Pyroprocessing uses electric potentials to 
pull actinides (uranium and 
\glspl{TRU}) onto a cathode from a molten salt bath, while undesired materials
either stay in the molten salt bath or on an anode \hl{Citation}

Reprocessing \gls{SNF} is known to improve reosurce utilization, compared 
with a once-through fuel cycle. Ine estimate reports that uranium ore
requirements would decrease by 25\% annually if \gls{SNF} is reprocessed 
to create \gls{MOX} fuel for \glspl{LWR} \cite{widder_benefits_2010}. 
Although this decrease in uranium ore requirements may improve the 
sustainability of the nuclear fuel cycle, current estimates of uranium 
reserves are expected to last for another 50-100 years 
\cite{widder_benefits_2010}. Therefore, this metric is only of concern 
if there is an increase in the price of uranium ore. Another known benefit 
of reprocessing is the reduction in the volume and radioacitivity of 
waste sent to a final repository. The volume of waste can be reduced by 
a factor of four by reprocessing \gls{SNF}, reducing the number of 
final reprositories needed to store all waste \cite{widder_benefits_2010}. 
The reduction of the radioacitivity of waste in a repository occurrs because 
of the removal of actinides, which are the long-lived nuclides that 
dominate radioactivity after the first 50 years of decay. 

Despite these benefits of reprocessing, it is not without its disdavantages.
The amount of \gls{LLW} generated greatly increases by using 
reprocessing \cite{widder_benefits_2010}. The \gls{LLW} generated 
from reprocessing includes any solvents or materials used for the process.
Pyroprocessing is expected to reduce the amount of \gls{LLW} comapred 
with aqueous reprocessing, because the molten salt bath can be cleaned 
and re-used. Therefore, the reprocessing technolgy employed affects 
this metric. Additionally, reporcessing is associated with an increase 
risk of proliferation. The increased risk for aqueous reprocessing 
comes from the creation of a stream of pure plutonium 
\cite{widder_benefits_2010}. This concern has led to the development of new 
methods to reprocess \gls{SNF} without creating pure plutonium, such 
as NUEX and COEX reprocessing methods \cite{widder_benefits_2010}. 
Additionally, reprocesisng \gls{SNF} is expected to cost more than 
using a once-through fuel cycle \cite{widder_benefits_2010}. Estimated 
\gls{LCOE} are \$30-75/MWh for a once-through fuel cycle, with an 
average of about \$47/MWh, and \$35-81/MWh for a fuel 
cycle with reprocessing, with an average of about \$52/MWh 
\cite{widder_benefits_2010}. The ranges in the \gls{LCOE} comes from 
differences in assumptions in determining the costs. 

\subsection{Uranium enrichment}
In the nuclear fuel cycle, enrichment increases the relative abundance of 
uranium-235 compared to uranium-238 in the fuel. This step is needed if the 
reactors deployed require a larger concentration of fissile material in 
the fuel that what is naturally ocurring. The capacity of an 
enrichment facility is based on a number of quantities \cite{tsoulfanidis_nuclear_2013}:
\begin{subequations}
    \begin{equation}
        F = \text{mass (kg) of feed material per unit time}
    \end{equation}
    \begin{equation}
        P = \text{mass (kg) of product material per unit time}
    \end{equation}
    \begin{equation}
        T = \text{mass (kg) of tails produced per unit time}
    \end{equation}
    \begin{equation}
        x_f = \text{weight fraction of $^{235}$U in the feed stream}
    \end{equation}
    \begin{equation}
        x_p = \text{weight fraction of $^{235}$U in the product stream}
    \end{equation}
    \begin{equation}
        x_t = \text{weight fraction of $^{235}$U in the tails stream}
    \end{equation}
    \begin{equation}
        SWU = \text{separtive work units (SWU); the physical work required to separate the isotopes}
    \end{equation}
\end{subequations}

Each of these quantites can be related using the following equations:
\begin{subequations}
    \begin{equation}
        F = P + T
    \end{equation}
    \begin{equation}
        x_fF = x_pP + x_tT
        \label{eq:enrichment_2}
    \end{equation}
    \begin{equation}
        SWU = [P*V(x_p) +T*V(x_t) - F*V(x_f)]*t
    \end{equation}
    \text{in which:}
    \begin{equation}
        V(x_i) = (2x_i - 1)*\ln\left(\frac{x_i}{1-x_i}\right)
    \end{equation}
    \begin{equation}
        t = \text{a given time period}
    \end{equation}
    \label{eq:enrichment}
\end{subequations}

The equations in Eq. \ref{eq:enrichment} show that as the desired enrichment level 
of the product feed increases, the feed mass must also increase, assuming everything 
else is held constant. The equations also show that as the desired product 
enrichemnt level increases, the amount of \gls{SWU} capacity needed also increases. 
However, if the product mass decreases, then the required \gls{SWU} capacity 
also decreases. Therefore, changes to the product enrichment level and mass 
throughput have opposite effects on the amount of feed material and required 
\gls{SWU} capacity. Changes to both quantities that are in opposite directions 
(i.e. an increase in one and a decrease in the other) will not lead to a clear 
expectation to the change in feed material of \gls{SWU} capacity. The effect 
on these variables will depend on the magnitude of the two changes.

\section{Research Goals}
The goal of this work is to investigate the impacts of deploying reactors fueled by 
\gls{HALEU} in the United States, including the impacts on the nuclear fuel cycle and 
on the reactors deployed. This goal is accomplished 
through multiple steps. The first step is modeling of multiple transition scenarios
to select advanced reactors requiring \gls{HALEU} fuel. The next is performing 
sensitivity analysis on the modeled transition scenarios and optimize the scenarios
based on perscribed metrics. Finally, the impact of the \gls{HALEU} production 
methods on reactor performance is investigated. 

The structure of this thesis is as follows. Chapter 2 provides some 
background information on the nuclear fuel cycle, nuclear fuel cycle 
modeling efforts, and the use of \gls{HALEU}-fuel in reactors.
Chapter 3 discusses the methodology for each section of this work, 
including the fuel cycle scenarios modeled, the advanced reactor selection 
method, and the neutronics models. Chapter 4 discusses the results of the 
once-through fuel cycle transition scenarios. Chapter 5 discusses the 
results of the recycling transition scenarios. Chapter 6 discusses the results 
of the sensitivity analysis and the optimization schemes applied to the 
transition scenarios. Chapter 7 discusses the impacts of various \gls{HALEU}
production methods on the performance of the reactors. Finally, Chapter 
8 provides some concluding remarks and suggestions for future work. 